\documentclass{article}
\usepackage[T1]{fontenc}
\usepackage{lmodern}
\usepackage{graphicx}
\usepackage{amsmath,amssymb}
\usepackage[a4paper,margin=2cm]{geometry}
\usepackage[absolute,overlay]{textpos}
\setlength{\TPHorizModule}{1mm}
\setlength{\TPVertModule}{1mm}
\usepackage[indonesian]{babel}
\usepackage{hyperref}
\setlength{\emergencystretch}{2em}
% unit: Halaman 36

\begin{document}

% === Halaman 36 (python/native) ===

• XOR Rcon: XOR-kan hasilnya dengan round constant (Rcon) ke i/4 (atau Rcon[i/4]).

-

Nilai Rcon berbeda-beda untuk setiap j = i/4, yaitu Rcon[j] = (RC]j], 0, 0, 0), 

dengan RC[1]=1, RC[j] = 2•RC[j-1], simbol • menyatakan perkalian yang 

didefinisikan di dalam GF(28). 

-

Nilai RC[j] di dalam heksadesimal adalah [STA11]: RC[1]=01, RC[2]=02, 

RC[3]=04, RC[4]=08, RC[5]=10, RC[6]=20, RC[7]=40, RC[8]=80, RC[9]=1B, 

RC[10]=36. 

• Simpan hasil fungsi g ke dalam peubah temp

c)

XOR-kan w[i-4] dengan temp

\end{document}
